\begin{thebibliography}{99}
\bibitem{[1]} 厚生労働省，``第3部　生活習慣調査の結果,'' 令和元年国民健康・栄養調査報告，pp.186，2019.
\bibitem{[2]} 大見拓寛, ``運転者の居眠り状態評価の画像センサ,'' 人工字臓器42巻1号，pp.99-103，2013.
\bibitem{[3]} 高木　和，中山　治人，岡田　志麻，牧川　方昭，``頭部の動きに着目した運転時における覚醒度低下検出の可能性,'' 生体医工学，pp.46-51，2013.
\bibitem{[4]} 金子成彦，藤田悦則，``ドライバーの覚醒低下警告・防止に向けた技術開発,'' 特集長時間運転と疲労，pp.57-63，2013.
\bibitem{[5]} Takumi MIYAZAWA，Ichiro FUKUMOTO，``生理指標を用いた居眠り防止システムの基礎検討,'' Journal of Advanced Science，Vol.20，No.1$\&$2，2008.
\bibitem{[6]} 横山直樹，``居眠り検知を目的とした生理学的パラメータの研究,'' ライフサポート，Vol.23，No.1，2011.
\bibitem{[7]} 谷田　陽介，萩原　啓，``心拍RRIのローレンツプロット情報に着目した入眠移行気の簡易推定法,'' 生体医工学，156-162，2006.
\bibitem{[8]} 豊福　史，山口　和彦，萩原　啓，``心電図RR間隔のローレンツプロットによる副交感神経活動の簡易推定法の開発,'' 人間工学，Vol.43，No.4，2007.
\bibitem{[9]} 松本　佳昭，森　信彰，三田尻　涼，江　鐘偉，``心拍揺らぎによる精神的ストレス評価法に関する研究,'' ライフサポート，Vol.22，No.3，2010.
\bibitem{[10]} 日本自律神経学会，自律神経機能検査　第三版，文光堂，2000.
\bibitem{[11]} 早野　純一郎，岡田　暁宣，安間　文彦，``心拍のゆらぎ: そのメカニズムと意義,'' 人工臓器25巻5号，pp.870-880，1996.
\end{thebibliography}